% SPDX-FileCopyrightText: Copyright (c) 2026 NVIDIA CORPORATION & AFFILIATES. All rights reserved.
% SPDX-License-Identifier: Apache-2.0
%
% ROLE: map every reader-facing name the paper uses for a registered decision, literature lesson,
%   stage, run or table label to the identifier the campaign record uses (restructure plan, section 6
%   row N; integrator decision 4). This appendix and appendix-reproducibility.tex are the only places
%   where amendment and lesson identifiers appear in reader-facing text.
% EVIDENCE: CR/CONTRACT.md (amendment headings A1-A21 and their numbered items),
%   CR/literature/LESSONS.md (lesson headings LR-01 to LR-15), CR/facts/tierA.json and
%   CR/facts/tierBC.json (runs and tuning rounds), CR/report/scripts/build_report_data.py (report row
%   labels), scripts/make_tables.py RELABEL (the paper's relabeling of those rows).
%   Public-content rule (A18): name no host, cluster, node, partition, account, job or internal URL.

\section{Names used in this paper and the campaign record}
\label{app:provenance}

The study's campaign record (\cref{app:repro}) identifies its registered decisions, literature
lessons, stages and runs by short labels. The paper uses descriptive names instead.
\Cref{tab:provenance} maps each name to the record's label, so a reader can find the registration,
lesson or run behind a statement. Amendments are the dated decisions recorded in the record's
contract during the study; lessons are the ranked findings of the study's literature review. Rows whose last column is empty are not referred to in the text and are
listed so that the numbering is complete.

{\small
\setlength{\LTcapwidth}{\linewidth}
\begin{longtable}{@{}P{0.42\linewidth}P{0.27\linewidth}P{0.24\linewidth}@{}}
\caption{Names used in this paper and the campaign record's identifiers. Amendments (A$n$) are in
  the record's contract, lessons (LR-$nn$) in its literature review; a decimal marks an item of an
  amendment.}\label{tab:provenance}\\
\toprule
Name in this paper & Record identifier & Where \\
\midrule
\endfirsthead
\toprule
Name in this paper & Record identifier & Where \\
\midrule
\endhead
\bottomrule
\endfoot
\multicolumn{3}{@{}l}{\textit{Registered decisions (amendments to the study's contract)}} \\*
The replicate rule (paired workload replicates replace identical re-runs, which are deterministic)
  and the pilot's two-clause noise rule & A1 & \cref{sec:replicates,sec:noise-floor,app:protocol} \\
The comparison against the validation-selected best baseline & A2.1 & \cref{sec:eval} \\
The service objective in this study (an ITL bound and an end-to-end slowdown bound against
  $\Ezero$; no TTFT bound) & A2.2 & \cref{sec:objective} \\
Windowed goodput & A2.3 & \cref{sec:objective} \\
Baselines evaluated as implemented, by design, with a brief sanity check & A2.5 &
  \cref{sec:baselines,app:ports} \\
AgentX replay through the study's own lowering, with play recycling & A3 & \cref{sec:agentx} \\
Duplicated trace copies & A4 & \cref{app:workloads} \\
Router-state lag as a robustness check, re-run on the test cells & A5.1 & \cref{sec:robustness} \\
Worker sets differ only in their count & A5.2 & \cref{sec:setup} \\
How fixes found during the study reach the upstream repositories & A6 & \\
Device-tier KV overlap passed to catalog policies in replay & A7 & \cref{tab:caveats,app:repro} \\
The repository mirror of the campaign record & A8 & \cref{app:repro} \\
The remote CPU lane and its cross-host parity & A9, A10 & \cref{app:parity} \\
Informed starts for M1's restarts & A11.1 & \cref{sec:restarts,app:starts} \\
The concentration audit, and the sign constraint registered in advance as its remedy & A11.2 &
  \cref{sec:gaming,sec:sign-constraint} \\
$\Nw = 6$ in both validation and test, reported as selection-exposed & A11.3 &
  \cref{sec:splits,sec:res-n} \\
Breadth-first scheduling of the main tuning on the remote lane; the lag patch in a separate build
  & A11.4, A11.5 & \cref{app:repro,app:results} \\
M1 restricted to its default-router start (a secondary test finalist) & A12 (``M1-default-init'')
  & \cref{sec:secondary,sec:res-ladder} \\
The live GPU check & A13 & \cref{sec:live} \\
The alternative SLO definitions & A14 & \cref{sec:robustness,sec:res-robustness} \\
M1-noaff (session coefficient pinned at 0) & A15 & \cref{sec:ladder} \\
The secondary \ais{}-informed league & A16 & \cref{app:ais-league} \\
The extended-feature arm M1-v2 & A17.1 & \cref{sec:ladder} \\
The headline-arm rule (the learned arm with the best score on fresh validation replicates, the rule
  that also picks the best baseline) & A17.2 & \cref{sec:selection} \\
The order in which arms would have been cut had compute run short & A17.3 & \\
The public-content rule (no internal infrastructure names in public material) & A18 & \\
Token-weighted goodput (secondary) & A19 & \cref{sec:a19} \\
The live finalists, registered before any finalist live job & A20 & \cref{sec:live-plan} \\
The \policy{lmetric} port fixed after the study froze its results & A21 & \cref{sec:ports}; the
  footnote at each headline statement \\
\midrule
\multicolumn{3}{@{}l}{\textit{Literature lessons}} \\*
Windowed goodput scored by a clipped mean log-ratio & LR-01 & \cref{sec:objective} \\
A noise floor from replicates and workload segments, not from repeats & LR-02 &
  \cref{sec:replicates,sec:noise-floor} \\
Common random numbers within each CMA-ES generation & LR-03 & \cref{sec:noise} \\
Paper-faithful baselines (not adopted: the baselines are evaluated as implemented) & LR-04 &
  \cref{sec:baselines} \\
Search spaces without flat directions; a suggested band of 10--30\% of decisions changed by the
  initial step; named context sources & LR-05 & \cref{sec:space,app:spaces} \\
Context and set-relative features that stay stable across $\Nw$ & LR-06 & \cref{sec:features} \\
The strongest heuristics inside M1's class, and restarts started from them & LR-07 &
  \cref{sec:features,sec:restarts} \\
SLO re-scoring at several scales (one TTFT threshold per family fails across 1K--128K prompts) &
  LR-08 & \cref{sec:objective,sec:robustness} \\
Sessions released causally; strata by load mode & LR-09 & \cref{sec:loadmodes,sec:strata} \\
Selection on validation over restarts, the same for baselines & LR-10 & \cref{sec:selection} \\
Segments as the unit of the pre-registered headline test; the intersection-union test; the
  probability-of-improvement criterion; performance profiles & LR-11 &
  \cref{sec:headline,fig:profiles} \\
Calibration per worker count, and the cache-pressure analysis & LR-12 &
  \cref{sec:calibration,fig:cache-pressure} \\
The worker-share cap and the other concentration flags & LR-13 & \cref{sec:gaming,tab:lr13} \\
The reading rule for the robustness checks (staleness and timing) & LR-14 & \cref{sec:robustness}
  \\
The drift check before M2, warm starts, and nest-level features for M3 & LR-15 &
  \cref{sec:drift,sec:nested} \\
\midrule
\multicolumn{3}{@{}l}{\textit{Stages}} \\*
The pilot & the pilot stage (\file{facts/pilot.json}) & \cref{sec:pilot} \\
The pilot's review & the gate (\file{facts/gate.json}) & \cref{sec:pilot,app:audits} \\
The main tuning & phase 2 & \cref{sec:training} \\
The first round of tuning (the ported baselines, M0, and M1 with its sign-constrained re-run) &
  tier A (\file{facts/tierA.json}) & \cref{sec:budget} \\
Later rounds: M1 continued, M1-noaff, M1-v2 and M2 rank 2; both arms of the queue-threshold
  ablation; the \ais{}-informed league & tier B; tier C; tier AIS (\file{facts/tierBC.json}) &
  \cref{sec:budget} \\
The first, unconstrained M1 & the tier-A M1 & \cref{sec:gaming} \\
The audit and fix midway through tuning & the phase-2 midpoint (\code{phase2-mid}) &
  \cref{sec:gaming,app:audits} \\
The single test pass & the Select+Test stage (``test stage'') & \cref{sec:eval} \\
The audit after the test pass (a completeness critic and three refuters) & the final audit
  (\file{audits/final/}) & \cref{sec:audits,app:audits} \\
The live GPU check & the live phase (A13, A20) & \cref{sec:live} \\
\midrule
\multicolumn{3}{@{}l}{\textit{Tools, cell set and runs}} \\*
The frozen cell set & \code{lr-cells-v5} & \cref{sec:freeze,app:repro} \\
The tuning harness; the evaluation and report tools & \code{lr-train}; \code{lr-eval},
  \code{lr-report} (harness version \code{lrh-4}) & \cref{sec:training,app:repro} \\
``restart $r$, gen.\ $g$, CMA-ES mean'' or ``\ldots, best-so-far sample'' & run key
  \code{<run>-s<r> g<g> mean} or \code{best\_so\_far}, kept in a \code{\% run key:} comment on each
  row of the generated tables & \cref{tab:leaderboard-full} \\
Run prefixes of the learned arms & \code{m1c} (M1), \code{m1} (M1 unconstrained), \code{m1cont},
  \code{m1noaff}, \code{m1v2}, \code{m2r2}, \code{m0} (M0) & \\
Run prefixes of the ablation and the \ais{}-informed league & \code{ablam1}, \code{ablabase};
  \code{m0ais}, \code{m1ais}, \code{m2ais}, \code{llmdobm} & \\
\midrule
\multicolumn{3}{@{}l}{\textit{Row labels in tables and figures}} \\*
M1 & ``M1 (= A12 M1-default-init)'' & \cref{tab:ladder-results,fig:leaderboard} \\
M1 unconstrained (flagged for gaming) & ``M1 unconstrained (gaming-flagged)'', ``(tier A,
  gaming-flagged)'' & \cref{tab:ladder-results,tab:lr13} \\
M1-noaff; M1-v2 & ``M1-noaff (A15)''; ``M1-v2 (A17.1)'' & \cref{tab:ladder-results} \\
M1 or ramjet + queue threshold (ablation) & ``(abl.\ a)'', ``(ablation a)'' &
  \cref{tab:ladder-results,fig:leaderboard} \\
SLO: ITL bound only; SLO: slowdown bound only; SLO: ITL + length-scaled TTFT; SLO: ITL + absolute
  E2E bound & ``A14 itl\_only'', ``A14 e2e\_only'', ``A14 ttft\_len\_itl'', ``A14 abs\_e2e\_itl'' &
  \cref{tab:robustness,fig:robustness} \\
Token-weighted goodput (secondary) & ``A19 good tokens (secondary)'' &
  \cref{tab:robustness,fig:robustness} \\
Worker with no good request & ``dump worker'' & \cref{tab:lr13} \\
round-robin; \policy{lmetric}; sticky-session bounded@defaults, sticky-session hard@defaults;
  llm-d-optimized-baseline@defaults (throughput) & ``round\_robin''; ``lmetric (port)'';
  ``sticky-bounded@defaults'', ``sticky-hard@defaults''; ``llm-d-optimized-baseline@throughput'' &
  \cref{fig:leaderboard,tab:leaderboard-full} \\
Test segments, for example Mooncake w3, AgentX T1, FAST25 synth.\ x0, sessions s6 &
  \code{mooncake:w3}, \code{agentx:T1}, \code{fast25\_synthetic:x0}, \code{sessions:s6} &
  \cref{fig:headline} \\
\end{longtable}}
